\section{Graph Characterization of Partial Views}\label{sec:graph-characterization}
%==============================================================================


Before defining the general confusability graph, we develop a concrete example: partial observations can induce nontrivial failure topology.

\subsection{A Running Example: The Binary Square}\label{sec:binary-square}

Consider a system with two binary facts $(x_1, x_2) \in \{0,1\}^2$, giving four latent states:
\[
(0,0), \quad (0,1), \quad (1,0), \quad (1,1).
\]
Suppose the architecture exposes two partial views:
\begin{itemize}
\item View 1 reveals only coordinate $x_1$
\item View 2 reveals only coordinate $x_2$
\end{itemize}

\paragraph{Confusability structure.}
Two states are confusable when at least one view fails to separate them:
\begin{itemize}
\item Through View 1: states $(0,0)$ and $(0,1)$ are confusable (both have $x_1=0$), and states $(1,0)$ and $(1,1)$ are confusable (both have $x_1=1$).
\item Through View 2: states $(0,0)$ and $(1,0)$ are confusable (both have $x_2=0$), and states $(0,1)$ and $(1,1)$ are confusable (both have $x_2=1$).
\end{itemize}
The confusability graph is therefore a 4-cycle:
\[
(0,0) \sim (0,1) \sim (1,1) \sim (1,0) \sim (0,0),
\]
where opposite corners $(0,0)$ and $(1,1)$, as well as $(0,1)$ and $(1,0)$, are \emph{not} adjacent. This is not a clique: the architecture can distinguish some pairs even though it cannot distinguish all pairs.

\paragraph{Exact recovery.}
Because the graph is 2-colorable (e.g., color by parity $c(x_1,x_2) = x_1 \oplus x_2$), exact recovery is possible with a 2-ary auxiliary tag. However, the graph is not 1-colorable (it contains edges), so a 1-tag decoder must fail on at least one confusable pair. Under the uniform source, the best 1-tag exact success probability is $1/2$ (choose either diagonal as the success set).

\paragraph{What this example shows.}
\begin{enumerate}
\item Partial views induce \emph{structured} confusability, not uniform ambiguity.
\item The confusability graph captures exactly which state pairs the architecture cannot separate.
\item Graph-theoretic properties (chromatic number, independence number) determine exact recovery costs.
\end{enumerate}

Figure~\ref{fig:failure-topology} contrasts this structured failure with the clique-shaped failure that arises when no observation provides any discrimination. The theorems below generalize this phenomenon to arbitrary multi-fact partial-view systems.

\subsection{General Graph Characterization}

This section turns the one-shot obstruction into the paper's main structural object: the confusability graph induced by admissible partial views. The single-fact converse family treats ambiguity classes that are effectively clique-shaped under the observation transcript. One extension is to let the latent state be a tuple of facts and let each location reveal only a subset of coordinates. Failure then becomes structured: some latent pairs remain indistinguishable and others do not, so the resulting confusability graph need not be complete.

In this extension, ambiguity is now over full latent tuples rather than single fact values: two tuples are confusable when the allowed partial observations fail to separate them. The induced graph is therefore a failure map for the architecture, not merely a bookkeeping device.

The graph can be handled either explicitly or implicitly. If there are $d$ facts over an alphabet of size $q$, then the latent state space has cardinality $q^d$, so any explicit confusability graph materialization already starts from an exponentially large vertex set.\leanmeta{\LH{MFT113}} A naive explicit construction ranges over at most $q^{2d}$ ordered state pairs.\leanmeta{\LHrng{MFT}{116}{117}} For that reason the appropriate algorithmic interface is implicit: the formalization packages adjacency as a Boolean confusability oracle on state pairs, proves that this oracle is equivalent to the confusability relation itself, and also packages an equivalent agreement-set oracle together with a linear bound on the total view-membership work needed by the direct scan implementation.\leanmeta{\LHrng{MFT}{114}{115}, \LHrng{MFT}{132}{135}} Under the standard explicit representation in which a latent tuple is stored as a length-$d$ coordinate array and each admissible view is stored as its coordinate subset, this one-shot oracle is therefore computable by direct inspection in deterministic time $O(d + \sum_{\ell}|V_\ell|)$, hence $O(Ld)$ for $L$ views: compute the agreement set of the two tuples once, then scan the views and test whether some view is contained in that agreement set. The theory below is therefore compatible with implicit graph access even when explicit graph materialization is impractical, while still leaving the one-shot adjacency test itself polynomial in the size of its explicit input.

The deterministic restriction is also not fundamental at the graph-construction level. The Lean development formalizes a support-overlap extension in which each location may emit any observation from a finite support set, and two states are confusable when some admissible location has overlapping observation support on those states. The present deterministic partial-view model embeds into that broader framework as the singleton-support special case, so the graph construction itself already extends beyond exact deterministic views even though the paper develops only the deterministic architectural arc.\leanmeta{\LHrng{MFT}{110}{112}}

At the same time, the exact current coordinate-subset model is \emph{not} graph-universal on the full tuple space, and the induced graph class can be characterized more sharply than mere vertex transitivity. The formalization first shows that equality at a view is equivalent to containment of that view inside the coordinate-agreement set of the two tuples, so confusability depends only on which coordinates agree, not on the actual symbols carried there.\leanmeta{\LHrng{MFT}{118}{119}, \LH{MFT122}} From any view family one therefore obtains an upward-closed family of coordinate subsets,
\[
\mathcal U := \{S \subseteq [F] : \exists \ell\ \text{with}\ V_\ell \subseteq S\},
\]
such that two distinct tuples are adjacent if and only if their agreement set lies in $\mathcal U$. Conversely, every upward-closed family $\mathcal U$ arises from some coordinate-view architecture, so on the full labeled tuple space the realizable confusability relations are exactly the agreement-set relations cut out by upward-closed families.\leanmeta{\LHrng{MFT}{125}{128}} When the alphabet has at least two symbols, every proper coordinate subset occurs as the agreement set of some distinct pair of tuples, so this characterization is complete up to the irrelevant full-agreement case corresponding to identical tuples.\leanmeta{\LHrng{MFT}{129}{131}} Equivalently, coordinatewise alphabet permutations preserve the confusability relation and induce graph automorphisms; for any two latent tuples there exists such an automorphism carrying one to the other.\leanmeta{\LHrng{MFT}{120}{124}} Thus the exact model realizes a specific monotone agreement-set graph class rather than an arbitrary finite graph family. In particular, non-vertex-transitive graphs such as a three-vertex path cannot arise as full confusability graphs in this exact model. What remains open is a classification up to unlabeled graph isomorphism, and the corresponding classification for restricted realized-state subsets or for the stochastic support-overlap extension.

\begin{theorem}[Exact recovery as graph colorability]\label{thm:multifact-colorability}
For a finite multi-fact latent tuple observed through a finite family of partial-coordinate views, exact zero-error recovery with a $T$-ary auxiliary tag exists if and only if the induced confusability graph is $T$-colorable.
\end{theorem}
\leanmeta{\LHrng{MFT}{3}{4}, \LH{MFT20}}

\begin{proof}
Let $G_{\mathrm{conf}}$ be the confusability graph.

\emph{Necessity.} Suppose exact recovery with $T$ tags is possible. If two latent tuples are adjacent in $G_{\mathrm{conf}}$, then by definition some allowed observation leaves them indistinguishable. Exact recovery therefore cannot assign them the same tag, because the joint observation-tag pair would then coincide on two distinct latent tuples and Theorem~\ref{thm:pair-injective} would be violated. Hence the tag assignment is a proper $T$-coloring of $G_{\mathrm{conf}}$.

\emph{Sufficiency.} Suppose $c$ is a proper $T$-coloring of $G_{\mathrm{conf}}$. Define the auxiliary tag of a latent tuple to be its color. If two latent tuples produce the same observation and the same color, then they are confusable and adjacent, contradicting proper colorability. Therefore the joint observation-color map is injective, and Theorem~\ref{thm:pair-injective} yields an exact decoder.
\end{proof}

An edge of $G_{\mathrm{conf}}$ records a specific pair of latent states that the view architecture cannot separate exactly, and a proper coloring is exactly a disambiguation budget that resolves those failures. Exact recovery cost is therefore determined by failure topology, not merely by the fact that more than one source exists.

This statement is packaged through an explicit simple confusability graph object, and for $n$ repeated copies the full block confusability graph is identified with the $n$-fold strong power of the one-shot confusability graph. The zero-error criterion is therefore literally a graph-colorability theorem on an explicit power family rather than only an analogy.\leanmeta{\LHrng{MFT}{34}{35}}

\begin{theorem}[Success-set characterization]\label{thm:multifact-success-set}
In the same model, exact decoding on a designated success set $S$ is equivalent to $T$-colorability of the induced confusability subgraph on $S$.
\end{theorem}
\leanmeta{\LH{MFT9}}

\begin{proof}
Restrict the argument of Theorem~\ref{thm:multifact-colorability} to the success set. Exactness is required only on $S$, so the coloring condition is required only on confusable pairs inside $S$.
\end{proof}

\begin{theorem}[Exact finite-error success budget]\label{thm:multifact-max-success}
For each tag alphabet size $T>0$, there is a largest success-set size
\[
M_T := \max\{|S| : S \text{ induces a } T\text{-colorable confusability subgraph}\}.
\]
An exact decoder with $T$ tags exists on a success set $S$ if and only if $|S|\le M_T$ after replacing $S$ by some $T$-colorable success set of size $M_T$. Under the uniform source, the optimal exact success probability is therefore $M_T/|\mathcal X|$.
\end{theorem}
\leanmeta{\LHrng{MFT}{14}{15}}

\begin{proof}
Theorem~\ref{thm:multifact-success-set} identifies exact success sets with induced $T$-colorable subgraphs. Since the latent alphabet is finite, a maximum-cardinality such set exists. The mechanized characterization exposes this optimum directly. In the binary four-cycle witness, the optimum at $T=1$ is exactly two states, so the best uniform exact success probability is $2/4=1/2$.
\end{proof}

The same characterization is also available in weighted form. For any finite source weight $w$ on the latent alphabet, the maximum mass of a $T$-colorable induced subgraph is the exact finite success value under budget $T$. This is the weighted version of the same structural claim: under limited budget, the best exact decoder is the largest colorable safe subworld allowed by the induced failure graph.\leanmeta{\LHrng{MFT}{18}{19}}

Equivalently, this optimum can be packaged as an explicit exact finite rate-distortion value for the extended model at tag budget $T$: every exact success set has weighted mass at most this value, and some exact decoder attains it. Thus the extended model no longer has only a converse inequality; at the one-shot level it has an exact weighted success law rather than merely an upper bound.\leanmeta{\LHrng{MFT}{23}{24}}

\begin{theorem}[First non-clique witness]\label{thm:multifact-nonclique}
There exists a binary two-fact system with single-coordinate observation views whose confusability graph is not a clique. In that system, two tags suffice for exact recovery, one tag is impossible, and any one-tag decoder can be exact on at most half of the four latent states under the uniform source.
\end{theorem}
\leanmeta{\LHrng{MFT}{5}{8}, \LH{MFT10}}

\begin{proof}
The witness is the four-state binary square with one location revealing the first coordinate and the other revealing the second. States that agree on one coordinate are confusable through the corresponding location, but opposite corners are not confusable, so the graph is a $4$-cycle rather than a clique. An explicit exact two-tag decoder is obtained by the parity coloring $c(x_1,x_2)=x_1\oplus x_2$. With one tag, any exact success set must be an independent set of the cycle, hence has size at most two, yielding success probability at most $1/2$ under the uniform source.
\end{proof}

\begin{theorem}[Block-composition law]\label{thm:multifact-product}
If two partial-observation systems are colorable with tag alphabets of sizes $T_1$ and $T_2$, then their block-composed product system is colorable with a tag alphabet of size $T_1T_2$.
\end{theorem}
\leanmeta{\LHrng{MFT}{11}{13}, \LH{MFT16}}

For later use, recall the strong product definition: if $G$ and $H$ are graphs, then $G\boxtimes H$ has vertex set $V(G)\times V(H)$, and $(u,v)$ is adjacent to $(u',v')$ exactly when either $u=u'$ and $v\sim v'$, or $u\sim u'$ and $v=v'$, or $u\sim u'$ and $v\sim v'$. The block-composition theorem is model-specific because the admissible product views generate exactly these three adjacency cases.

\begin{proof}
Take proper colorings of the two component confusability graphs and encode the pair of colors as one color in the product alphabet. If two product states are confusable, then some allowed product view fails to separate them. In the product system this means that either the first coordinates agree and the second coordinates are confusable, or the second coordinates agree and the first coordinates are confusable, or both coordinates are confusable. These are exactly the three adjacency cases in the strong product $G_1\boxtimes G_2$. Hence every confusable product pair is separated by at least one component coloring, and therefore by the paired color as well.
\end{proof}

This product law now has both directions in the formal development. The upper direction is the multiplicative coloring construction above. More sharply, once both location alphabets are nonempty, the product confusability graph is exactly the strong product of the two component confusability graphs. The graph object is therefore not just analogous to zero-error graph products; it is literally a strong-product construction for the induced failure map. \leanmeta{\LHrng{MFT}{25}{26}}

The lower direction is clique-based: if the first component contains a confusability clique of size $c_1$ and the second contains one of size $c_2$, then the product system contains a confusability clique of size $c_1c_2$, so any exact product decoder requires at least $c_1c_2$ tags. The same mechanism also propagates to finite-error success sets: the largest exactly decodable success set in the product system is at least the product of the largest exactly decodable success sets in the two components. At the one-shot zero-error coding level, the maximum independent-set code size also obeys the same product lower bound. Thus the extended model now has an honest strong-product graph law together with matching lower/upper budget consequences.\leanmeta{\LH{MFT17}}

At this point the induced graph class has four explicit structural features. It contains non-clique graphs, as witnessed by the binary square. It contains the cluster-collapse subclass on which the later equality theorem is exact. It is closed under the block-composition operation, which becomes strong product on confusability graphs. And, by the agreement-set characterization above, it is exactly a monotone agreement-set class rather than an arbitrary graph family on the full tuple space. The iterated family below packages these features into a concrete infinite class.

\paragraph{Example family: iterated binary squares.}
The four-state witness above is only the smallest member of a larger model-generated family. Let $W$ denote the binary square system of Theorem~\ref{thm:multifact-nonclique}. Composing $m$ independent copies of $W$ yields a system with $2m$ binary fact coordinates and $4^m$ latent states whose confusability graph is the $m$-fold strong power
\[
C_4^{\boxtimes m}.
\]
Because each copy of $W$ is exactly recoverable with two tags, repeated application of Theorem~\ref{thm:multifact-product} gives exact recovery on the $m$-fold product with $2^m$ tags. Conversely, each copy contains a confusability clique of size $2$, so the clique lower bound in the same product theorem yields a confusability clique of size $2^m$ in the product system. Hence every exact decoder for this family requires at least $2^m$ tags, and therefore the exact tag budget is
\[
T_{\mathrm{exact}}(W^{\otimes m}) = 2^m.
\]
This provides an infinite family of non-clique examples, generated directly by the model, on which the graph object, the exact budget, and the strong-product law all interact nontrivially rather than only in the single four-cycle case. At this point the paper has moved decisively beyond the trivial threshold story: above the unit-rate corner, not all failures are equally severe, because the architecture can induce a non-clique topology with a strictly more structured recovery law. \leanmeta{\LHrng{MFT}{11}{13}, \LHrng{MFT}{16}{17}, \LHrng{MFT}{29}{30}}

Once exact one-shot recovery is governed by graph colorability, repeated composition naturally lifts the theory to strong graph powers and asymptotic zero-error rates. The next subsection makes that lift explicit.
